\section{Longest Palindromic Subsequence}
\label{sec:dp-longest-palindromic-subsequence}

\problemstatement{Given a string $s$, find the length of its
\textbf{longest palindromic subsequence} (a subsequence that reads the
same forward and backward).}

\begin{patternbox}
\emph{``Longest subsequence that is a palindrome''} reduces directly to
Section~\ref{sec:dp-lcs}: a palindrome reads identically forward and
backward, so a subsequence of $s$ is a palindrome exactly when it is
\emph{also} a subsequence of $s$'s \textbf{reverse} -- making the longest
palindromic subsequence of $s$ exactly the longest common subsequence of
$s$ and $\textit{reverse}(s)$.
\end{patternbox}

\subsection*{Understanding the reduction}

Let $r = \textit{reverse}(s)$. Any common subsequence of $s$ and $r$
is, by construction, a sequence of characters appearing in order in $s$
that \emph{also} appears in order in $r$ -- but appearing in order in $r$
is the same as appearing in \emph{reverse} order in $s$. A sequence that
is both a subsequence of $s$ in forward order and a subsequence of $s$
in reverse order is precisely a palindromic subsequence of $s$. So:
\[
  |\textit{LPS}(s)| = |\textit{LCS}(s, \textit{reverse}(s))|.
\]

\subsection*{All Four Approaches}

Every stage is Section~\ref{sec:dp-lcs}'s LCS implementation, called
with $s_1 = s$ and $s_2 = \textit{reverse}(s)$.

\begin{lstlisting}
#include <bits/stdc++.h>
using namespace std;

int lcsSpaceOptimized(string& s1, string& s2) {
    int n = s1.size(), m = s2.size();
    vector<int> prev(m + 1, 0), current(m + 1, 0);

    for (int i = 1; i <= n; i++) {
        for (int j = 1; j <= m; j++) {
            if (s1[i - 1] == s2[j - 1]) {
                current[j] = 1 + prev[j - 1];
            } else {
                current[j] = max(prev[j], current[j - 1]);
            }
        }
        prev = current;
    }
    return prev[m];
}

int longestPalindromicSubsequence(string s) {
    string r = s;
    reverse(r.begin(), r.end());
    return lcsSpaceOptimized(s, r);
}

int main() {
    cout << longestPalindromicSubsequence("bbbab") << endl; // 4
    return 0;
}
\end{lstlisting}

\subsection*{Dry run}

\dryrun{Take $s = \texttt{"bbbab"}$.}

$r = \texttt{"babbb"}$. The LCS of \texttt{"bbbab"} and \texttt{"babbb"}
is \texttt{"bbbb"} (length $4$), which is indeed a palindrome and a
subsequence of the original \texttt{"bbbab"} (drop the \texttt{'a'}).

\begin{complexitybox}
\textbf{Time Complexity.} $O(n^2)$ (since $|s_1|=|s_2|=n$), inherited
directly from LCS.
\textbf{Space Complexity.} $O(n)$ for the rolling-row LCS.
\end{complexitybox}

\begin{takeawaybox}
``Longest palindromic X'' problems frequently reduce to ``longest common
X between the string and its reverse'' -- a reduction worth checking
first before writing a dedicated palindrome recurrence (which does also
exist for this specific problem, expanding outward from the center, but
the LCS reduction shown here requires no new derivation at all).
\end{takeawaybox}
